% SOURCE_STATEMENT: "Is the instability of anti-de-Sitter space caused by resonant mode mixing that transfers energy from low to high frequencies until an apparent horizon forms?"
% Preparation and final self-audit: CONTRIBUTING.md prompts C and D; source/audit artifacts retained in the local physics_statements_2026_10_11 batch.
% !TeX program = lualatex
\documentclass[11pt]{article}
\usepackage[margin=25mm]{geometry}
\usepackage{fontspec}
\setmainfont{Latin Modern Roman}
\usepackage{amsmath,amssymb,mathtools}
\usepackage{unicode-math}
\setmathfont{Latin Modern Math}
\usepackage{xurl,hyperref}
\hypersetup{colorlinks=true,urlcolor=blue}
\setcounter{secnumdepth}{0}
\setlength{\parindent}{0pt}
\setlength{\parskip}{5pt}
\setlength{\emergencystretch}{3em}
\newcommand{\sref}[2]{\hyperlink{source:#1:#2}{[S#2]}}
\newcommand{\cataloguescope}[1]{\paragraph{Recorded target.}#1}
\newcommand{\R}{\mathbb R}
\newcommand{\C}{\mathbb C}
\newcommand{\N}{\mathbb N}
\newcommand{\Z}{\mathbb Z}
\newcommand{\Q}{\mathbb Q}
\newcommand{\F}{\mathbb F}
\newcommand{\eps}{\varepsilon}
\begin{document}
% UnsolvedMath ID 30002150; source catalogue code OWR-12009-007
\setcounter{section}{30002149}
\section{Resonant Mode Mixing in Anti-de-Sitter Instability}\label{30002150}
{\small UnsolvedMath ID 30002150 · Mathematical Physics}
\subsection{Definitions and mathematical statement}

\paragraph{Shared notation.}$\mathbb N=\{1,2,\ldots\}$; the standard number fields have their usual meanings. All additional conventions are stated locally.


\paragraph{Definitions and conventions.}
Anti-de Sitter spacetime is the maximally symmetric Lorentzian solution with negative cosmological constant. Its linearized field modes have frequencies; resonant interactions occur when combinations of those frequencies match. A transfer of energy to increasing frequencies is a turbulent cascade. An apparent horizon is a marginally outer trapped surface on a spacelike slice; its formation is a geometric diagnostic, distinct from a proof of linear mode instability.

\paragraph{Target retained from the canonical source catalogue.}
\paragraph{Catalogue wording retained for traceability.}
Is the instability of anti-de-Sitter space caused by resonant mode mixing that transfers energy from low to high frequencies until an apparent horizon forms?

\paragraph{Source scope and qualification.}
This asks for a mechanism in the source evolution problem. Resonance in a truncated mode system is not by itself proof that a horizon forms in the full field equations.
\subsection{Short English statement}
Does resonant transfer of energy to shorter scales explain black-hole formation in anti-de Sitter instability?
\subsection{Sources}
\begin{itemize}
\item[\hypertarget{source:30002150:1}{S1}] ULAM.AI and the UnsolvedMath Contributors, UnsolvedMath, record 30002150 (OWR-12009-007). Catalogue curation under CC BY 4.0; definitions and English exposition added locally.
\url{https://huggingface.co/datasets/ulamai/UnsolvedMath}.
\item[\hypertarget{source:30002150:2}{S2}] Mathematical Aspects of General Relativity (2012), Oberwolfach Reports, DOI 10.4171/owr/2012/37, pp. 2269--2333. Primary report accessed 2026-10-11; relevant PDF pages 52.
\url{https://doi.org/10.4171/owr/2012/37}.
\item[\hypertarget{source:30002150:3}{S3}] Bizon--Rostworowski, weak turbulence (catalogue-listed context; not independently consulted in this preparation).
\url{https://doi.org/10.1103/PhysRevLett.107.031102}.
\item[\hypertarget{source:30002150:4}{S4}] Bizon--Maliborski--Rostworowski, resonant dynamics (catalogue-listed context; not independently consulted in this preparation).
\url{https://doi.org/10.1103/PhysRevLett.115.081103}.
\end{itemize}
\label{end:30002150}
\end{document}
